\clearpage\SourceLeaf{463}
\IndexHeading{NUMERISCHES RECHNEN.}
\IndexSubheading{Anzeige.}
\IndexRow{0}{\textsc{Leonelli}, Logarithmische Supplemente}{.\,.\,.\,.\,.\,.\,.\,.\,.\,.\,.\,.\,.\,.\,.\,.\,.}{Seite}{121}
\IndexSubheading{Nachlass.}
\IndexRow{0}{Vorschriften, um den Logarithmen des Sinus eines kleinen Bogens zu finden}{.\,.\,.\,.\,.}{—}{128}
\IndexRow{0}{Interpolation der Cotangenten und Cosecanten kleiner Bögen}{.\,.\,.\,.\,.\,.\,.\,.\,.\,.\,.}{—}{129}
\IndexRow{0}{Musterrechnung, um aus $A=p\cos P$, $B=p\sin P$ $p$ und $P$ zu finden}{.\,.\,.\,.\,.\,.\,.}{—}{130}
\IndexHeading{WAHRSCHEINLICHKEITSRECHNUNG.}
\IndexSubheading{Nachlass und Briefwechsel.}
\IndexRow{0}{Zwei Aufgaben aus der Wahrscheinlichkeitsrechnung}{.\,.\,.\,.\,.\,.\,.\,.\,.\,.\,.\,.\,.\,.}{—}{133}
\IndexRow{0}{Zur Geschichte der Entdeckung der Methode der kleinsten Quadrate}{.\,.\,.\,.\,.\,.\,.\,.}{—}{136}
\IndexRow{0}{Kritische Bemerkungen zur Methode der kleinsten Quadrate}{.\,.\,.\,.\,.\,.\,.\,.\,.}{—}{142}
\IndexRow{0}{Kleinere Beiträge zur Methode der kleinsten Quadrate}{.\,.\,.\,.\,.\,.\,.\,.\,.\,.\,.\,.\,.}{—}{149}
\IndexRow{0}{Eine Ausgleichsformel für Mortalitätstafeln}{.\,.\,.\,.\,.\,.\,.\,.\,.\,.\,.\,.\,.\,.\,.}{—}{155}
\IndexHeading{GRUNDLAGEN DER GEOMETRIE.}
\IndexSubheading{Anzeigen, Nachlass und Briefwechsel.}
\IndexRow{0}{Über die ersten Gründe der Geometrie\quad \textsc{Gauss} an \textsc{Bolyai} 1799 Dec. 16.}{.\,.\,.\,.\,.\,.}{—}{159}
\IndexRow{\widthof{Über die ersten Gründe der Geometrie\quad}}{\textsc{Gauss} an \textsc{Bolyai} 1804 Nov. 25.}{.\,.\,.\,.\,.\,.}{—}{160}
\IndexRow{0}{Einige Sätze die ersten Gründe der Geometrie betreffend (\textit{Nachlass})}{.\,.\,.\,.\,.\,.\,.\,.\,.}{—}{163}
\IndexRow{0}{Zur Theorie der Parallellinien.\quad \textsc{Gauss} an \textsc{Olbers} 1806 Juli 30.}{.\,.\,.\,.\,.\,.\,.\,.\,.\,.}{—}{165}
\IndexRow{\widthof{Zur Theorie der Parallellinien.\quad}}{Aus \textsc{Schumachers} Tagebuch 1808 Nov.}{.\,.\,.\,.\,.\,.}{—}{165}
\IndexRow{\widthof{Zur Theorie der Parallellinien.\quad}}{Ideen (\textit{Nachlass} 1813)}{.\,.\,.\,.\,.\,.\,.\,.\,.\,.\,.\,.\,.}{—}{166}
\IndexRow{0}{\textsc{Legendres} Theorie der Parallelen:\quad \textsc{Gerling} an \textsc{Gauss} 1816 März 11.}{.\,.\,.\,.\,.\,.\,.}{—}{167}
\IndexRow{\widthof{\textsc{Legendres} Theorie der Parallelen:\quad}}{\textsc{Gauss} an \textsc{Gerling} 1816 April 11.}{.\,.\,.\,.\,.\,.\,.}{—}{168}
\IndexRow{0}{\textsc{Schwab}, Commentatio in primum elementorum Euclidis librum; \textsc{Metternich}, Vollständige Theorie der Parallellinien (\textit{Anzeige} 1816 April 20.)}{.\,.\,.\,.\,.\,.\,.\,.\,.\,.\,.\,.\,.}{—}{170}
\IndexRow{0}{Die transcendente Trigonometrie.\quad \textsc{Wachter} an \textsc{Gauss} 1816 Dec. 12.}{.\,.\,.\,.\,.\,.\,.\,.}{—}{175}
\IndexRow{\widthof{Die transcendente Trigonometrie.\quad}}{\textsc{Gauss} an \textsc{Olbers} 1817 April 28.}{.\,.\,.\,.\,.\,.\,.\,.}{—}{177}
\IndexRow{0}{Astralgeometrie:\quad \textsc{Gerling} an \textsc{Gauss} 1818 Juli 23.}{.\,.\,.\,.\,.\,.\,.\,.\,.\,.\,.\,.}{—}{178}
\IndexRow{\widthof{Astralgeometrie:\quad}}{\textsc{Gauss} an \textsc{Gerling} 1818 Aug. 25.}{.\,.\,.\,.\,.\,.\,.\,.\,.\,.\,.\,.\,.\,.}{—}{179}
\IndexRow{\widthof{Astralgeometrie:\quad}}{\textsc{Gerling} an \textsc{Gauss} 1819 Jan. 25.}{.\,.\,.\,.\,.\,.\,.\,.\,.\,.\,.\,.}{—}{179}
\IndexRow{\widthof{Astralgeometrie:\quad}}{\textsc{Gauss} an \textsc{Gerling} 1819 März 16.}{.\,.\,.\,.\,.\,.\,.\,.\,.\,.\,.\,.\,.}{—}{181}
\IndexRow{0}{\textsc{Müller}, Theorie der Parallelen (\textit{Anzeige} 1822 Oct. 28.)}{.\,.\,.\,.\,.\,.\,.\,.\,.\,.\,.\,.\,.}{—}{183}
\IndexRow{0}{Zur Parallelentheorie:\quad \textsc{Gauss} an \textsc{Taurinus} 1824 Nov. 8.}{.\,.\,.\,.\,.\,.\,.\,.\,.\,.\,.\,.}{—}{186}
\IndexRow{\widthof{Zur Parallelentheorie:\quad}}{\textsc{Gauss} an \textsc{Olbers} 1827 Mai 3.}{.\,.\,.\,.\,.\,.\,.\,.\,.\,.}{—}{188}
\IndexRow{\widthof{Zur Parallelentheorie:\quad}}{\textsc{Gauss} an \textsc{Schumacher} 1827 Oct. 11.}{.\,.\,.\,.\,.\,.\,.\,.\,.\,.\,.}{—}{189}
\IndexRow{0}{Über die Winkel des Dreiecks (\textit{Nachlass})}{.\,.\,.\,.\,.\,.\,.\,.\,.\,.\,.\,.\,.\,.\,.\,.\,.}{—}{190}
\IndexRow{0}{Zur Theorie der geraden Linie und der Ebene (\textit{Nachlass})}{.\,.\,.\,.\,.\,.\,.\,.\,.\,.\,.\,.}{—}{193}
\IndexRow{0}{Über die ersten Gründe der Geometrie:\quad \textsc{Gauss} an \textsc{Bessel} 1829 Jan. 27.}{.\,.\,.\,.\,.\,.}{—}{200}
\IndexRow{\widthof{Über die ersten Gründe der Geometrie:\quad}}{\textsc{Bessel} an \textsc{Gauss} 1829 Febr. 10.}{.\,.\,.\,.\,.\,.}{—}{201}
\IndexRow{\widthof{Über die ersten Gründe der Geometrie:\quad}}{\textsc{Gauss} an \textsc{Bessel} 1830 April 9.}{.\,.\,.\,.\,.\,.\,.}{—}{201}
\IndexRow{0}{Zur Theorie der Parallellinien (\textit{Nachlass} 1831)}{.\,.\,.\,.\,.\,.\,.\,.\,.\,.\,.\,.\,.\,.}{—}{202}
